\chapter{Maximal Extractable Value}\label{ch:m3:maximal-extractable-value}

A trader sends a transaction that buys a million dollars of ether from a public pool, accepting up to
half a percent of slippage. Before it is included, someone else's purchase of ether lands in the same
block just ahead of it, and a sale of that ether just behind it. The trader receives the least his
tolerance allowed; the other party pockets the difference, less what it paid the builder of the block for
the privilege of those two positions. Nothing in the exchange was hacked. The value came from the right to
order transactions, which on a \omterm{def:m3:blockchains-for-traders:chain}{blockchain} belongs to whoever builds the block, and which is bought and
sold in a market of its own. This chapter explains that market: where extractable value comes from, the
sandwich with its arithmetic, the arbitrages and liquidations that are its benign cousins, the auctions and
the separation of roles through which blocks are now built, and the private channels and designs meant to
keep traders' orders out of it.

\section{Ordering power and extractable value}

\begin{definition}[Maximal extractable value, searcher]\label{def:m3:maximal-extractable-value:mev}
\emph{Maximal extractable value}\index{maximal extractable value} (MEV) is the value that can be obtained by
including, excluding or reordering transactions in a block beyond the ordinary fees. A
\emph{searcher}\index{searcher} is a participant, usually an automated system, that watches pending
transactions and the chain's state for such opportunities and submits transactions to capture them.
\end{definition}

The term was introduced as ``miner extractable value'' by Daian and co-authors, who documented bots that bid up
transaction fees against each other in open auctions to be first in line for arbitrage opportunities on
\omterm{def:m3:automated-market-makers:amm}{decentralised exchanges}, likened them to high-frequency traders, and warned that the fees paid for priority posed a
risk to the chain's consensus. The value has three sources. A transaction may create a price discrepancy that
someone can trade on (a large swap moves a pool). A change of external prices may create one (an \omterm{def:m3:automated-market-makers:amm}{automated market
maker}'s stale price, \cref{ch:m3:automated-market-makers}). And a protocol may pay whoever acts first (a
liquidation bonus, \cref{ch:m3:defi-credit-and-leverage}). In each case the value goes to whoever can be first, or
around, in the order of a block.

\begin{definition}[Front-running, back-running]\label{def:m3:maximal-extractable-value:front}
\emph{Front-running}\index{front-running} is placing a transaction ahead of a known pending one to profit from the
price move that the pending one will cause. \emph{Back-running}\index{back-running} is placing a transaction
immediately after a known pending one, to capture the opportunity it leaves (an arbitrage, a liquidation) without
harming it.
\end{definition}

\section{Front-running and the sandwich}

\begin{definition}[Sandwich attack, slippage tolerance]\label{def:m3:maximal-extractable-value:sandwich}
A \emph{sandwich attack}\index{sandwich attack} is a front-run and a back-run around a victim's swap in the same pool:
the attacker buys what the victim is buying, the victim buys at the higher price, and the attacker sells after it.
The victim's \emph{slippage tolerance}\index{slippage tolerance} is the worst price, expressed as a percentage below the
quote at submission, at which it allows its swap to execute; below it, the swap reverts.
\end{definition}

\begin{omfigure}
\centering
\begin{tikzpicture}[font=\small,
  t/.style={draw,thick,rounded corners=2pt,align=center,minimum height=0.85cm,text width=3.0cm},
  arr/.style={-{Stealth[length=2.2mm]},thick}]
\node[t,draw=omExo,fill=omExo!8] (a) at (0,0) {1.\ attacker buys ETH\\(front-run)};
\node[t,draw=omThm,fill=omThm!8] (v) at (4,0) {2.\ victim buys ETH\\at its worst price};
\node[t,draw=omExo,fill=omExo!8] (b) at (8,0) {3.\ attacker sells ETH\\(back-run)};
\draw[arr] (a) -- (v); \draw[arr] (v) -- (b);
\node[font=\footnotesize,anchor=north,align=center,text width=11cm] at (4,-0.8) {one block, one pool: the builder places the three
transactions consecutively; the victim's tolerance bounds how far step 1 may push the price};
\end{tikzpicture}

\omcaption{A \omterm{def:m3:maximal-extractable-value:sandwich}{sandwich attack}. The attacker's two trades bracket the victim's in the same block; the victim's \omterm{def:m3:maximal-extractable-value:sandwich}{slippage tolerance} is
the only limit on the size of the first. Schematic.}\label{fig:m3:maximal-extractable-value:sandwich}
\end{omfigure}

\begin{proposition}[The tolerance bounds the sandwich]\label{prop:m3:maximal-extractable-value:bound}
Let a victim swap $\Delta x$ of \omterm{def:m3:blockchains-for-traders:key}{token} A for \omterm{def:m3:blockchains-for-traders:key}{token} B in a constant-product pool with minimum output $(1 - \tau)$ times the
quote at submission. The victim's output falls as the attacker's front-run grows, so there is a largest front-run
$a^*(\tau)$ that leaves the victim at exactly its minimum; any larger front-run makes the victim's swap revert and the
attack fail. The victim's loss is at most $\tau$ times its quoted output, and $a^*$ and the attacker's gross gain grow with
$\tau$. With $\tau = 0$ no sandwich is possible.
\end{proposition}

\begin{proof}
A front-run moves the pool's price of B up, so the victim's output (\cref{prop:m3:automated-market-makers:swap}) is decreasing
in the front-run; the feasible set is an interval $[0, a^*]$, with $a^*$ increasing in $\tau$. The victim receives at least its
minimum, so it loses at most $\tau$ of its quote. At $\tau = 0$ only a zero front-run is feasible.
\end{proof}

Within the feasible interval the attacker's gain rises with the size of its front-run, so it attacks at the bound: the tutorial
finds the bound by binary search, in the pool's integer arithmetic. For a USD~1 million purchase of ether in a pool holding
USD~15 million and 5{,}000 ether with a 0.30\% fee, a tolerance of 0.5\% allows a front-run of USD~38{,}912 and a gross gain of
USD~5{,}067; the victim receives 1.56 ether fewer than its quote, about USD~4{,}700 at the pool's starting price. At 1\%, the
front-run doubles to USD~78{,}119, the gain to USD~10{,}130 and the victim's loss to 3.12 ether.

\begin{omfigure}
\begin{tikzpicture}
\begin{axis}[width=11cm,height=5cm,
  xlabel={victim's slippage tolerance, \%},ylabel={USD thousand},
  tick label style={font=\small},label style={font=\small},
  xmin=0,xmax=2,ymin=0,ymax=22,grid=major,grid style={gray!25},
  legend columns=2,legend style={font=\footnotesize,at={(0.5,-0.3)},anchor=north,draw=none,fill=none,/tikz/every even column/.append style={column sep=8pt}},
  table/col sep=comma]
\addplot[omExo,very thick] table[x=tol_pct,y=gross_k]{figdata/markets-3/22-maximal-extractable-value/tolerance.csv};\addlegendentry{attacker's gross gain}
\addplot[omThm,very thick,dashed] table[x=tol_pct,y=loss_k]{figdata/markets-3/22-maximal-extractable-value/tolerance.csv};\addlegendentry{victim's loss (at 3{,}000 per ether)}
\end{axis}
\end{tikzpicture}

\omcaption{The sandwich of a USD~1 million purchase of ether in a pool of USD~15 million and 5{,}000 ether (0.30\% fee), at the largest
front-run the victim's tolerance allows (\cref{prop:m3:maximal-extractable-value:bound}). Both grow linearly with the tolerance. Integer
pool arithmetic; illustrative pool. Data: the chapter's tutorial.}\label{fig:m3:maximal-extractable-value:tolerance}
\end{omfigure}

\begin{omfigure}
\begin{tikzpicture}
\begin{axis}[width=11cm,height=5cm,
  xlabel={front-run, USD thousand},ylabel={attacker's gross gain, USD thousand},
  tick label style={font=\small},label style={font=\small},
  xmin=0,xmax=78,ymin=0,ymax=10.5,grid=major,grid style={gray!25},unbounded coords=jump,
  legend columns=2,legend style={font=\footnotesize,at={(0.5,-0.3)},anchor=north,draw=none,fill=none,/tikz/every even column/.append style={column sep=8pt}},
  table/col sep=comma]
\addplot[gray,thick,dashed] table[x=front_k,y=gross_k]{figdata/markets-3/22-maximal-extractable-value/curve.csv};\addlegendentry{if the victim did not revert}
\addplot[omExo,very thick] table[x=front_k,y=feasible_k]{figdata/markets-3/22-maximal-extractable-value/curve.csv};\addlegendentry{feasible at 0.5\% tolerance}
\end{axis}
\end{tikzpicture}

\omcaption{The attacker's gross gain against the size of its front-run for the swap of \cref{fig:m3:maximal-extractable-value:tolerance}
at a 0.5\% tolerance. The gain keeps rising, but beyond USD~38{,}912 the victim's swap would revert and the attacker would be left holding
ether bought at a premium. Data: the chapter's tutorial.}\label{fig:m3:maximal-extractable-value:curve}
\end{omfigure}

The arithmetic tells a trader what to do. Tolerance is an option written to anyone who can order the block; it should be set from the
expected price impact of the trade itself plus a small margin, not a round number; large swaps should be split or sent through channels
where the order is not public (below).

\section{Back-running, CEX--DEX arbitrage and liquidation races}

Most extracted value does not harm a particular user.

\begin{definition}[CEX--DEX arbitrage]\label{def:m3:maximal-extractable-value:cexdex}
\emph{CEX--DEX arbitrage}\index{CEX--DEX arbitrage} is trading an \omterm{def:m3:automated-market-makers:amm}{automated market maker}'s stale price back to the price on \omterm{def:m3:centralised-exchanges:cex}{centralised
exchanges} after they move, usually hedged at once on a \omterm{def:m3:centralised-exchanges:cex}{centralised exchange}, and placed first in the next block.
\end{definition}

It is the \omterm{def:m3:automated-market-makers:lvr}{loss-versus-rebalancing} of \cref{ch:m3:automated-market-makers} seen from the arbitrageur's side. In the tutorial's pool, a 1\%
rise of ether on \omterm{def:m3:centralised-exchanges:cex}{centralised exchanges} lets an arbitrageur buy about USD~52{,}000 of ether from the pool with a profit of about USD~180
after the pool's fee; a 2\% rise, about USD~127{,}000 with a profit of about USD~1{,}070. The profit is a race: whoever's transaction is first
in the block takes it all, and the competition is over the right to be first. Liquidation races are the same race for a lending protocol's
liquidation bonus. \omterm{def:m3:maximal-extractable-value:front}{Back-running} a large swap to arbitrage the pool it has moved is harmless to the swapper, and the private channels below
return part of its value to the swapper.

\section{Gas auctions, proposer--builder separation and bundles}

\begin{definition}[Priority gas auction]\label{def:m3:maximal-extractable-value:pga}
A \emph{priority gas auction}\index{priority gas auction} is the open competition in which \omterm{def:m3:maximal-extractable-value:mev}{searchers} bid up the \omterm{def:m3:blockchains-for-traders:gas}{priority fees} of transactions in
the public \omterm{def:m3:blockchains-for-traders:mempool}{mempool}, each outbidding the others, to be ordered first for the same opportunity.
\end{definition}

\omterm{def:m3:maximal-extractable-value:pga}{Priority gas auctions} spam the network with failed and replaced transactions and give the right to order to the block producer, who can also take
the opportunity itself. Since Ethereum's move to \omterm{def:m3:blockchains-for-traders:pow}{proof of stake} the market has been organised differently.

\begin{definition}[Proposer--builder separation, block builder, MEV relay, transaction bundle]\label{def:m3:maximal-extractable-value:pbs}
\emph{Proposer--builder separation}\index{proposer--builder separation} is the division of block production between \emph{block
builders}\index{block builder}, specialised parties that assemble the most valuable blocks from public transactions and private bundles and bid for
the right to have them proposed, and proposers (\omterm{def:m3:blockchains-for-traders:pow}{validators}) that choose the highest bid. An \emph{MEV relay}\index{MEV relay} is an intermediary that
receives builders' blocks, keeps their contents hidden from the proposer until it commits to propose, and forwards the best bid. A \emph{transaction
bundle}\index{transaction bundle} is an ordered list of transactions that a \omterm{def:m3:maximal-extractable-value:mev}{searcher} submits to builders, to be included together and in that order,
or not at all.
\end{definition}

\begin{omfigure}
\centering
\begin{tikzpicture}[font=\small,
  s/.style={draw,thick,rounded corners=2pt,align=center,minimum height=0.9cm,text width=2.2cm},
  arr/.style={-{Stealth[length=2.2mm]},thick}]
\node[s,draw=omDef,fill=omDef!8] (u) at (0,0) {users and\\wallets};
\node[s,draw=omExo,fill=omExo!8] (se) at (2.9,0) {searchers\\(bundles)};
\node[s,draw=omThm,fill=omThm!8] (bu) at (5.8,0) {block\\builders};
\node[s,draw=omProp,fill=omProp!8] (re) at (8.7,0) {MEV relays};
\node[s,draw=gray,fill=gray!8] (pr) at (11.6,0) {proposer\\(validator)};
\draw[arr] (u) -- (se); \draw[arr] (se) -- (bu); \draw[arr] (bu) -- (re); \draw[arr] (re) -- (pr);
\draw[arr,dashed] (u.south) to[bend right=18] node[below,font=\footnotesize] {private order flow} (bu.south);
\end{tikzpicture}

\omcaption{The supply chain of an Ethereum block under \omterm{def:m3:maximal-extractable-value:pbs}{proposer--builder separation}: \omterm{def:m3:maximal-extractable-value:mev}{searchers} send bundles to builders, builders send blocks and
bids to relays, and the proposer signs the most valuable block without seeing its contents first. \omterm{def:m3:maximal-extractable-value:private}{Private order flow} goes straight to builders.
Schematic.}\label{fig:m3:maximal-extractable-value:pbs}
\end{omfigure}

\begin{dated}{2026-09}{Proposer--builder separation in practice}\label{dat:m3:maximal-extractable-value:boost}
Flashbots describes MEV-Boost as open-source middleware run by \omterm{def:m3:blockchains-for-traders:pow}{validators} to access a competitive block-building market, its implementation of
\omterm{def:m3:maximal-extractable-value:pbs}{proposer--builder separation} for proof-of-stake Ethereum: builders send blocks to relays, relays forward the most profitable block to the proposer.
Testimony in a 2025 criminal trial put the share of \omterm{def:m3:blockchains-for-traders:pow}{validators} using MEV-Boost in April 2023 at 95\%. Relayscan's statistics for the 24 hours to 24
September 2026 counted 6{,}526 blocks delivered through relays, of about 7{,}200 slots, of which three builders built about 93\%; the
MEV-Boost dashboard maintained by Toni Wahrstätter put the share of MEV-Boost blocks at 91.58\% on the same day, with 7 relays and 25
builders active over fourteen days.
\end{dated}

Builders compete in an auction for each slot, and \omterm{def:m3:maximal-extractable-value:mev}{searchers} compete for position within builders' blocks by paying them: in a competitive market a
\omterm{def:m3:maximal-extractable-value:mev}{searcher} keeps only a small share of an opportunity's value and pays the rest to the builder, which passes most of it to the proposer. With a 90\%
payment, the sandwich of the previous section nets its \omterm{def:m3:maximal-extractable-value:mev}{searcher} about USD~487 of its USD~5{,}067, and the builder about USD~4{,}560. The value
extracted from the victim ends mostly with the \omterm{def:m3:blockchains-for-traders:pow}{validator}.

The system relies on trust in relays and their code. Prosecutors in New York charged two brothers with exploiting a flaw in MEV-Boost relay software in
April 2023: according to the government's post-trial filing, they baited other \omterm{def:m3:maximal-extractable-value:mev}{searchers}' sandwich bots with transactions, registered \omterm{def:m3:blockchains-for-traders:pow}{validators} to
obtain a block's private contents from a relay by submitting an invalid block header, and replaced the bundled transactions to take more than USD~26
million. The trial ended in a mistrial on 7 November 2025 when the jury failed to reach a verdict; the defence then moved for acquittal, a motion
argued on 23 March 2026.

\section{Private order flow, rollup sequencing and mitigations}

\begin{definition}[Private order flow, order-flow auction]\label{def:m3:maximal-extractable-value:private}
\emph{Private order flow}\index{private order flow} is transactions sent directly to builders or through private channels rather than to the public
\omterm{def:m3:blockchains-for-traders:mempool}{mempool}, where \omterm{def:m3:maximal-extractable-value:mev}{searchers} cannot see them. An \emph{order-flow auction}\index{order-flow auction} is a mechanism that lets \omterm{def:m3:maximal-extractable-value:mev}{searchers} bid for the right to
back-run a user's transaction, returning part of the proceeds to the user, without revealing the transaction in full.
\end{definition}

\begin{dated}{2026-09}{A private channel and an order-flow auction}\label{dat:m3:maximal-extractable-value:protect}
Flashbots Protect sends transactions to a private \omterm{def:m3:blockchains-for-traders:mempool}{mempool} hidden from \omterm{def:m3:maximal-extractable-value:front}{front-running} and sandwich bots, includes them only if they do not revert, and
refunds MEV and \omterm{def:m3:blockchains-for-traders:gas}{priority fees}. Its MEV-Share node shares selected information about a user's transaction with \omterm{def:m3:maximal-extractable-value:mev}{searchers}, simulates their bundles and
forwards successful ones to builders on condition that the user is paid a specified share of the MEV they create (by default 90\%); it accepts only
back-runs.
\end{dated}

Private channels cut sandwiches but move the trust: the user relies on the channel and the builders it sends to not to exploit the flow themselves,
and flow that goes to a few builders concentrates block building further, which is one reason why three builders build most blocks. On \omterm{def:m3:blockchains-for-traders:rollup}{rollups}, the
\omterm{def:m3:blockchains-for-traders:rollup}{sequencer} (\cref{ch:m3:blockchains-for-traders}) orders transactions; if it processes them first come, first served and keeps its queue private, there is
no public \omterm{def:m3:blockchains-for-traders:mempool}{mempool} to front-run, and the race moves to latency to the \omterm{def:m3:blockchains-for-traders:rollup}{sequencer}, as on a centralised venue. The designs of earlier chapters address the
same problem from the other side: \omterm{def:m3:automated-market-makers:routing}{intent-based routing} and batch auctions (\cref{ch:m3:automated-market-makers}) clear orders at one price, so that order
within a batch does not matter.

\section{Tutorial: anatomy of a sandwich}\label{tut:m3:maximal-extractable-value:sandwich}

\begin{tutorial}
\textbf{Goal.} Given a pool and a victim's swap with its tolerance, find the attacker's largest front-run, its profit after gas and the builder's
payment, and the victim's loss; show that the tolerance bounds both. \textbf{End state:}
\cref{fig:m3:maximal-extractable-value:tolerance,fig:m3:maximal-extractable-value:curve} and the numbers of the weekend problem.
\begin{tutsteps}
\item \textbf{The attack.} Front-run, victim, back-run, in the pool's integer arithmetic; binary search for the bound.
\omcode{code/firm/sandwich/firm_sandwich.py}{26}{61}{The sandwich at the largest feasible front-run.}\label{lst:m3:maximal-extractable-value:sandwich}
\item \textbf{Tolerances.} \texttt{by\_tolerance()} with and without a 90\% payment to the builder.
\item \textbf{Detection.} Find the pattern in a block's swaps: one sender on both sides of another's trade in the same pool.
\omcode{code/firm/sandwich/firm_sandwich.py}{64}{77}{Detecting a sandwich in a block's list of swaps.}\label{lst:m3:maximal-extractable-value:detect}
\item \textbf{Run} \texttt{cex\_dex\_arbitrage(3030)} and \texttt{fig\_mev.py}.
\end{tutsteps}
\textbf{What to change next.} Let the victim split its swap in ten across blocks; add a competing \omterm{def:m3:maximal-extractable-value:mev}{searcher} who bids for the same sandwich and find the
builder's share at which the attack stops paying; run \texttt{find\_sandwiches} on a synthetic block with two overlapping attacks.
\end{tutorial}

\section{Build: the sandwich simulator}\label{bld:m3:maximal-extractable-value:sandwich}

\begin{build}
\bfield{Purpose} The miniature firm swaps on chain as a customer and trades arbitrages as a \omterm{def:m3:maximal-extractable-value:mev}{searcher}: it must know what its own transactions expose,
set tolerances, price back-runs and detect when it has been sandwiched.
\bfield{Interface} \texttt{run(ra, rb, dx\_v, dx\_a)}; \texttt{min\_out(ra, rb, dx\_v, tolerance\_bps)}; \texttt{max\_front\_run}; \texttt{sandwich(ra,
rb, dx\_v, tolerance\_bps, gas\_cost, builder\_share\_bps)} returning the front-run, the victim's output and loss, and gross and net gains;
\texttt{find\_sandwiches(txs)}. Swaps from \texttt{firm.amm}.
\bfield{Rules} Integer base units and the pool's exact rounding; the victim's minimum computed from the quote at submission, as wallets do; the attack
refused if the victim would revert.
\bfield{Acceptance tests} \texttt{code/firm/sandwich/tests/}: the bound is exact to one unit; gain and loss grow with tolerance; a builder's share lowers
the net; zero tolerance leaves nothing; a sandwich detected in a block and not across pools.
\bfield{Stretch} Concentrated-liquidity pools crossing ticks; multi-pool sandwiches; back-run pricing for an \omterm{def:m3:maximal-extractable-value:private}{order-flow auction}.
\end{build}

\begin{omsources}
\item P.~Daian et al., ``Flash Boys 2.0: Frontrunning, Transaction Reordering, and Consensus Instability in Decentralized Exchanges'',
arXiv:1904.05234, 2019.
\item Flashbots documentation (MEV-Boost, MEV-Share, Flashbots Protect), GitHub repository flashbots/flashbots-docs, September 2026; relayscan.io
and mevboost.pics, 24 September 2026.
\item United States v.\ Peraire-Bueno, No.\ 1:24-cr-00293 (S.D.N.Y.): government's memorandum of 16 January 2026 (Doc 236) and the docket's
declaration of mistrial of 7 November 2025; joint motion for acquittal (Doc 225, 12 December 2025); order setting oral argument (Doc 240).
\end{omsources}

\section{Exercises}

\begin{exercise}[$\star$]\label{exo:m3:maximal-extractable-value:1}
Why can no one sandwich a swap sent with zero \omterm{def:m3:maximal-extractable-value:sandwich}{slippage tolerance}, and why do wallets not default to zero?
\end{exercise}

\begin{exercise}[$\star$]\label{exo:m3:maximal-extractable-value:2}
Classify as \omterm{def:m3:maximal-extractable-value:front}{front-running}, \omterm{def:m3:maximal-extractable-value:front}{back-running} or neither: a liquidation of an undercollateralised loan placed right after an \omterm{def:m3:blockchains-for-traders:contract}{oracle} update; an arbitrage
placed right after a large swap; a purchase placed right before a large purchase.
\end{exercise}

\begin{exercise}[$\star$]\label{exo:m3:maximal-extractable-value:3}
What does an \omterm{def:m3:maximal-extractable-value:pbs}{MEV relay} guarantee the builder, and what does it guarantee the proposer?
\end{exercise}

\begin{exercise}[$\star\star$]\label{exo:m3:maximal-extractable-value:4}
The victim's tolerance doubles from 0.5\% to 1\%. By how much do the attacker's gross gain and the victim's loss change?
\end{exercise}

\begin{exercise}[$\star\star$]\label{exo:m3:maximal-extractable-value:5}
A \omterm{def:m3:maximal-extractable-value:mev}{searcher} pays the builder 90\% of the gross gain and USD~20 of gas. What does it net on the 1\% sandwich?
\end{exercise}

\begin{exercise}[$\star\star$]\label{exo:m3:maximal-extractable-value:6}
Why does \omterm{def:m3:maximal-extractable-value:private}{private order flow} tend to concentrate block building?
\end{exercise}

\begin{exercise}[$\star\star\star$]\label{exo:m3:maximal-extractable-value:7}
\emph{Coding.} With \texttt{cex\_dex\_arbitrage}, find the smallest centralised price rise that makes the arbitrage worth USD~20 of gas.
\end{exercise}

\begin{exercise}[$\star\star\star$]\label{exo:m3:maximal-extractable-value:8}
\emph{Find the flaw.} ``We send our swaps through a private channel, so we pay no MEV.''
\end{exercise}

\section{Problem: Anatomy of a Sandwich}

\begin{problem}[{Weekend problem --- a USD 1 million swap and the searcher who saw it}]\label{pb:m3:maximal-extractable-value:1}
A trader buys ether with USD~1 million in a constant-product pool holding USD~15 million and 5{,}000 ether with a 0.30\% fee, through the public
\omterm{def:m3:blockchains-for-traders:mempool}{mempool}. A \omterm{def:m3:maximal-extractable-value:mev}{searcher}'s two transactions cost USD~20 of gas.

\medskip\noindent\textbf{Part I --- The quote.}
\begin{enumerate}
\item How much ether does the quote promise, and what is the swap's own price impact?
\item What is the trader's minimum output at 0.5\% tolerance? At 1\%?
\item What is the largest front-run at each tolerance?
\item Why is that the \omterm{def:m3:maximal-extractable-value:mev}{searcher}'s best choice?
\item What would the \omterm{def:m3:maximal-extractable-value:mev}{searcher} do at 0\% tolerance?
\end{enumerate}

\medskip\noindent\textbf{Part II --- The profit.}
\begin{enumerate}[resume]
\item What is the \omterm{def:m3:maximal-extractable-value:mev}{searcher}'s gross gain at each tolerance?
\item What does the trader lose, in ether and at 3{,}000 dollars per ether?
\item What does the \omterm{def:m3:maximal-extractable-value:mev}{searcher} net if it pays the builder nothing? If it pays 90\%?
\item Who ends up with most of the value?
\item Why does a competitive builder market push the \omterm{def:m3:maximal-extractable-value:mev}{searcher}'s share towards zero?
\end{enumerate}

\medskip\noindent\textbf{Part III --- Defences.}
\begin{enumerate}[resume]
\item What tolerance would you set for this swap, and why?
\item What would splitting the swap in ten across ten blocks change?
\item What would a private channel change, and what would it not?
\item How would an intent or a batch auction change the outcome?
\item How would you detect after the fact that the swap was sandwiched?
\end{enumerate}

\medskip\noindent\textbf{Part IV --- Judgement.}
\begin{enumerate}[resume]
\item Is a sandwich \omterm{def:m3:maximal-extractable-value:front}{front-running} in the sense of the market-abuse rules that bind brokers in regulated markets?
\item Why are back-runs and \omterm{def:m3:maximal-extractable-value:cexdex}{CEX--DEX arbitrage} treated differently from sandwiches?
\item What does the 2023 relay exploit show about the block-building supply chain?
\item State the \emph{named result}: the \omterm{def:m3:maximal-extractable-value:mev}{searcher}'s maximal net profit and the trader's loss at 0.5\% and 1\% tolerance.
\item In one sentence: what is a \omterm{def:m3:maximal-extractable-value:sandwich}{slippage tolerance}?
\end{enumerate}
\end{problem}

\section{Interview questions}

\begin{interviewq}[$\star$ \iqroles{trader}]\label{iq:m3:maximal-extractable-value:1}
What is a \omterm{def:m3:maximal-extractable-value:sandwich}{sandwich attack} and how do you protect your swaps from it?
\end{interviewq}

\begin{interviewq}[$\star$ \iqroles{developer}]\label{iq:m3:maximal-extractable-value:2}
Explain \omterm{def:m3:maximal-extractable-value:pbs}{proposer--builder separation} and the role of relays.
\end{interviewq}

\begin{interviewq}[$\star\star$ \iqroles{researcher}]\label{iq:m3:maximal-extractable-value:3}
Derive the largest front-run a victim's tolerance allows in a constant-product pool.
\end{interviewq}

\begin{interviewq}[$\star\star$ \iqroles{trader, researcher}]\label{iq:m3:maximal-extractable-value:4}
How is \omterm{def:m3:maximal-extractable-value:cexdex}{CEX--DEX arbitrage} related to a pool's \omterm{def:m3:automated-market-makers:lvr}{loss-versus-rebalancing}?
\end{interviewq}

\begin{interviewq}[$\star\star$ \iqroles{risk}]\label{iq:m3:maximal-extractable-value:5}
What risks does a firm running \omterm{def:m3:maximal-extractable-value:mev}{searchers} take, operationally and legally?
\end{interviewq}

\begin{interviewq}[$\star\star\star$ \iqroles{developer, researcher}]\label{iq:m3:maximal-extractable-value:6}
Design an \omterm{def:m3:maximal-extractable-value:private}{order-flow auction} that returns value to users without leaking their trades.
\end{interviewq}
